% ============================================================
%  part1/ch03-microscope.tex
%  Chapter 3: The Microscope
% ============================================================

\chapter{The Microscope}
\label{ch:microscope}

\epigraph{%
  To see a World in a Grain of Sand\\
  And a Heaven in a Wild Flower,\\
  Hold Infinity in the palm of your hand\\
  And Eternity in an hour.%
}{William Blake}

\section{The paradox of magnification}

The microscope reverses a natural intuition.

More magnification, one might think, means more information.
The finer the resolution, the closer we are to reality.
To see smaller things is to see more truly.

Yet this is not what happens.

At high magnification, context disappears.
A biological tissue seen at the cellular level reveals cell
structure but loses organ structure.
A metal surface seen at the atomic level reveals crystal lattice
geometry but loses the shape of the object.
A neural spike train resolved to the millisecond reveals
individual firing patterns but loses behavioural context.

Observation gains detail while losing organisation.

The microscope proves that information loss occurs not only by
coarse-graining, but also by over-resolution.

\section{Multiscale structure and scale selection}

Most real systems have meaningful structure at many scales
simultaneously.

\begin{definition}[Multiscale system]
  A \emph{multiscale system} is a system $X$ that can be
  decomposed as
  \[
    X = \bigcup_{\ell > 0} X_\ell,
  \]
  where $X_\ell$ denotes the structure of $X$ that is
  visible at scale $\ell$.
\end{definition}

Examples are ubiquitous.

\begin{example}[Scales of biological organisation]
  A living organism has structure at the scales of:
  atoms, molecules, organelles, cells, tissues, organs,
  organ systems, whole organisms, populations, and ecosystems.
  Observation at any single scale reveals some structure while
  hiding structure at all other scales.
\end{example}

A microscope selects a scale window:
\[
  \pi_{\ell_0} : X \to X_{\ell_0}.
\]
This is a projection.
Like the map and the photograph, it has a nontrivial kernel.

The context lost by magnification is
\[
  K_{\ell_0} = X \setminus X_{\ell_0}.
\]
Thus observation at small scale destroys access to large-scale
organisation, and observation at large scale destroys access to
fine structure.

\section{Admissibility depends on the reconstruction task}

The key insight of this chapter is that the right scale of
observation is not determined by the system alone.
It is determined by the reconstruction task.

Suppose a reconstruction task $R$ requires both local detail
$d$ and global context $c$.
Let $I(\ell)$ denote the information accessible at scale $\ell$.
Then neither extreme is optimal:
\[
  \lim_{\ell \to 0} I(\ell)
\]
captures microscopic detail while losing context, whereas
\[
  \lim_{\ell \to \infty} I(\ell)
\]
captures context while losing detail.

\begin{definition}[Admissible information for a task]
  For a reconstruction task $R$, define the
  \emph{admissible information} at scale $\ell$ as
  \[
    A_R(\ell) = \mathrm{Rel}(I(\ell), R),
  \]
  where $\mathrm{Rel}(\cdot, R)$ measures the relevance of the
  available information to the reconstruction goal $R$.
\end{definition}

The admissible scale is then
\[
  \ell^* = \arg\max_\ell A_R(\ell).
\]

The critical observation is that admissibility is not
maximised by information quantity:
\[
  \arg\max_\ell A_R(\ell)
  \neq
  \arg\max_\ell |I(\ell)|.
\]

\principle{
  Useful information is not maximum information.\\[4pt]
  Admissibility is always relative to a reconstruction goal.
}

This principle quietly anticipates every major framework developed
later in this book.

In CLIO (Part~V), admissibility is formalised as a constraint on
which projections are permitted.
In the admissibility sheaf (Chapter~\ref{ch:sheaf}),
admissible local sections are exactly those relevant to global
reconstruction.
In simulated agency (Chapter~\ref{ch:agency}), agents act only
through admissible perturbations of their local field.
In cosmology (Part~VII), the admissible observational data are
those that survive projection through causal horizons.

All of these are versions of the same recognition: the right
observation is not the most detailed one, but the one
appropriate to the task.

\section{Bidirectional information loss}

The map and photograph chapters illustrated information loss by
coarse-graining: the projection discards fine detail.

The microscope chapter introduces the opposite direction.
Over-resolution also loses information, by discarding large-scale
organisation.

\begin{proposition}[Bidirectional information loss]
  For a multiscale system $X = \bigcup_\ell X_\ell$, the
  total information accessible at any single scale $\ell_0$ is
  strictly less than the total information in $X$:
  \[
    |I(\ell_0)| < |I(X)|
    \qquad \text{for all } \ell_0.
  \]
  In particular, neither the limit $\ell_0 \to 0$ nor the limit
  $\ell_0 \to \infty$ recovers the full structure of $X$.
\end{proposition}

This means that complete information about a multiscale system is
not accessible from any single scale of observation.
Reconstruction of the full system requires integrating information
across multiple scales.

\begin{example}[Astronomy across scales]
  Astronomical observation operates simultaneously at scales
  ranging from individual photons to the large-scale structure
  of the universe.
  A spectrum reveals chemical composition.
  An image reveals spatial structure.
  A light curve reveals temporal variation.
  No single measurement captures all three.
  Reconstructing the physical state of a distant object requires
  combining observations made at different scales, resolutions,
  and wavelengths.
\end{example}

\section{Admissibility as a filter, not a threshold}

A common misconception is that admissibility means passing some
minimum quality threshold.
More data is assumed to be better, up to the point where
resolution is sufficient.

The microscope chapter corrects this.

Admissibility is not a threshold.
It is a filter that depends on the relationship between the
available data and the reconstruction goal.

Data can be inadmissible by being too coarse.
Data can also be inadmissible by being too fine.

An observation is admissible if and only if it captures the
structure relevant to the task while discarding the structure
irrelevant to it.

This bidirectional notion of admissibility will reappear in
Chapter~\ref{ch:sheaf} when we ask which local sections of the
admissibility sheaf can be glued into a global reconstruction,
and in Chapter~\ref{ch:open-problems} where we state as a
conjecture that the admissibility cohomology of cosmological
observer domains is nontrivial.

\section{Summary}

The map chapter introduced projection and the nontrivial kernel.

The photograph chapter introduced the admissibility class and the
reconstruction problem.

This chapter introduces the central principle that will govern the
rest of the book:

\begin{enumerate}[leftmargin=2em]
  \item Information loss is bidirectional: both coarse-graining
    and over-resolution destroy information.
  \item The right level of observation is determined by the
    reconstruction task, not by the system alone.
  \item Useful information is not maximum information.
    Admissibility is always relative to a goal.
\end{enumerate}

The reader who has absorbed these three points is already thinking
in the terms that the rest of this book will make precise.
